\documentclass[11pt]{article}

% --- Packages ---
\usepackage[utf8]{inputenc}
\usepackage[margin=1in]{geometry} % Better margins
\usepackage{amsmath, amssymb, amsthm} % Essential math tools
\usepackage{hyperref} % Clickable links/TOC
\usepackage{xcolor} % Custom colors
\usepackage{listings} % For code snippets
\usepackage{enumitem} % Better control over lists
\usepackage{ulem} % For underlining
\usepackage{enumitem} % to enumerate important items with bolded numbers
\usepackage{amsmath} % to highlight important equations with a box
\usepackage{empheq}

% --- Custom Math Environments ---
\newtheorem{theorem}{Theorem}[section]
\theoremstyle{definition}
\newtheorem{definition}{Definition}[section]
\newtheorem{example}{Example}[section]

% --- Code Snippet Styling ---
\lstset{
    basicstyle=\ttfamily\small,
    breaklines=true,
    keywordstyle=\color{blue},
    commentstyle=\color{gray},
    frame=single,
    numbers=left,
    numberstyle=\tiny\color{gray}
}

\setlength{\parindent}{0pt}      % Sets indentation to zero
\setlength{\parskip}{1em}        % Adds a gap between paragraphs

% --- Document Info ---
\title{CS 109: Intro to Probability \\ \large Lecture 1: Fundamentals of Probability}
\author{Chris Laganiere}
\date{\today}

\begin{document}

\maketitle
\tableofcontents
\newpage

\section{Definitions}
In probability, we often deal with sample spaces and events. 

\begin{definition}[Sample Space]
    The set of all possible outcomes of an experiment is called the \textbf{sample space}, denoted by $S$ or $\Omega$.
\end{definition}

\begin{definition}[Event]
    A subset of the sample space, with some specific outcomes of an experiment, denoted by $E$.
\end{definition}
    
\section{Fundamentals of Probability}
An event $E$ is a subset of a sample space $S$:

\begin{align}
    \{ E \subseteq S \}
\end{align}

Examples:

\begin{itemize}
    \item Tossing a coin: $S = \{H, T\}$
    \item 1 "Heads" on 2 flips: $S = \{[H, T], [H, H], [T, H]\}$
\end{itemize}

The measured probability of an event measured by the fraction of the time it occurs in experimentation.

\begin{align}
    P(E) = \lim_{n \to \infty} \frac{count(E)}{n}
\end{align}

Probability is always expressed from 0 to 1:

\begin{empheq}[box=\fbox]{align}
    0 \leq P(E) \leq 1
\end{empheq}

Equally likely outcomes are expressed with vertical pipes:

\begin{align}
    P(E) = \frac{\#event outcomes}{\# space outcomes} = \frac{|E|}{|S|}
\end{align}

\section{Axioms of Probability}

\begin{enumerate}[label=\textbf{\arabic*.}, ref=\arabic*]
    \item $0 \leq P(E) \leq 1$
    \item $P(S) = 1$
    \item if events $E$ and $F$ are mutually exclusive:
        \begin{empheq}[box=\fbox]{align}
            P(E \cup F) = P(E) + P(F)
        \end{empheq}
\end{enumerate}

It's sometimes easier to find probability of the complement of an event rather than the probability of the event itself.
\begin{align}
    P(E \cup E^C) &= 1 \\
                  &= P(E) + P(E^C)
\end{align}

The complement can be used to calculate the event's probability.

\begin{empheq}[box=\fbox]{align}
    P(E^C) &= 1 - P(E^C)
\end{empheq}

\end{document}